\NativeSlot{47}{57}{%
\ExamGroup{二、选择（将正确答案的序号填在括号里）。{\normalfont\kaishu\fontsize{10.8bp}{14bp}\selectfont （每题2分，共14分）}}

\par\noindent%
\ExamTile{.74}{%
% p102_s01

\ExamQuestion{1.}{5}%
万老师用三根木条做成一个等腰三角形衣架（如图），底边比腰长。这三根木条的长度可能是\mbox{（\hspace{1.8em}）}。\par
\ExamOptions{A. \(2.1\,\mathrm{dm}\)，\(2.1\,\mathrm{dm}\)，\(5.2\,\mathrm{dm}\)}{B. \(2\,\mathrm{dm}\)，\(3\,\mathrm{dm}\)，\(3\,\mathrm{dm}\)}{C. \(2.4\,\mathrm{dm}\)，\(2.4\,\mathrm{dm}\)，\(4.2\,\mathrm{dm}\)}{}
}\hfill%
\ExamTile{.235}{%
\BookFigure{figures/p102_s01-01.tex}
}\par
}
\NativeSlot{147}{90}{%
\par\noindent%
\ExamTile{.67}{%
% p102_s02

\ExamQuestion{2.}{5}%
下列说法正确的是\mbox{（\hspace{1.8em}）}。\par ①\(6.995\)精确到百分位是\(7.00\)。\par ②右图竖式的计算过程中运用了乘法分配律。\par ③\(0.2\)与\(0.4\)之间只有一个小数。\par ④相机支架一般都是三角形状，这是利用了三角形的稳定性。\par
\ExamOptions{A. ①②④}{B. ①③④}{C. ①④}{}
}\hfill%
\ExamTile{.305}{%
\BookFigure{figures/p102_s02-01.tex}
}\par
}
\NativeSlot{300}{45}{%
% p102_s03

\ExamQuestion{3.}{5}%
优优、佳佳、浩浩三名同学在一次跳绳比赛中的平均成绩是\(125\)个，乐乐跳了\(145\)个，他们四人跳的平均成绩是\mbox{（\hspace{1.8em}）}个。\par
\ExamOptions{A. 125}{B. 130}{C. 145}{}
}
\NativeSlot{376}{101}{%
% p102_s04

\ExamQuestion{4.}{5}%
下列的关系图中，能正确表示同一平面内两条直线位置关系的是\mbox{（\hspace{1.8em}）}。\par
\BookFigure{figures/p102_s04-01.tex}
}
\NativeSlot{545}{103}{%
% p102_s05

\ExamQuestion{5.}{5}%
李华在课中研究三角形，用直尺、圆规绘制了下面的三幅图，你知道他在研究什么问题吗？\mbox{（\hspace{1.8em}）}\par
\ExamOptions{A. 三角形的特征}{B. 三角形的内角和}{C. 三角形三边关系}{D. 三角形的周长}
\BookFigure{figures/p102_s05-01.tex}
}
\NativeSlot{719}{44}{%
% p102_s06

\ExamQuestion{6.}{5}%
短跑赛场上有两组队员正在准备比赛，采取三局两胜制。如果比赛中每个人都正常发挥，根据“田忌赛马”的策略，按照\mbox{（\hspace{1.8em}）}方案安排出场顺序，二组才能获胜。\par
}
